% \iffalse
%
% beamerthemeosx.dtx --- docTeX file for Beamer Theme OSX
%
% Copyright (C) 2026 Didier Verna
%
% Author: Didier Verna <didier@didierverna.net>
%
% This file is part of Beamer Theme OSX.
%
% Beamer Theme OSX may be distributed and/or modified under the conditions of
% the LaTeX Project Public License, either version 1.3c of this license or (at
% your option) any later version. The latest version of this license is in
% http://www.latex-project.org/lppl.txt and version 1.3c or later is part of
% all distributions of LaTeX version 2008 or later.
%
% Beamer Theme OSX consists of the files listed in the file `MANIFEST.md'.
%
%
% Commentary:
%
%
% Code:
%
%<*driver>
\documentclass{ltxdoc}
\usepackage[T1]{fontenc}
\usepackage{graphicx}
\EnableCrossrefs
\CodelineIndex
\RecordChanges
\begin{document}
  \DocInput{beamerthemeosx.dtx}
  \PrintChanges
  \PrintIndex
\end{document}
%</driver>
% \fi
%
% \CheckSum{0}
% \CharacterTable
%  {Upper-case    \A\B\C\D\E\F\G\H\I\J\K\L\M\N\O\P\Q\R\S\T\U\V\W\X\Y\Z
%   Lower-case    \a\b\c\d\e\f\g\h\i\j\k\l\m\n\o\p\q\r\s\t\u\v\w\x\y\z
%   Digits        \0\1\2\3\4\5\6\7\8\9
%   Exclamation   \!     Double quote  \"     Hash (number) \#
%   Dollar        \$     Percent       \%     Ampersand     \&
%   Acute accent  \'     Left paren    \(     Right paren   \)
%   Asterisk      \*     Plus          \+     Comma         \,
%   Minus         \-     Point         \.     Solidus       \/
%   Colon         \:     Semicolon     \;     Less than     \<
%   Equals        \=     Greater than  \>     Question mark \?
%   Commercial at \@     Left bracket  \[     Backslash     \\
%   Right bracket \]     Circumflex    \^     Underscore    \_
%   Grave accent  \`     Left brace    \{     Vertical bar  \|
%   Right brace   \}     Tilde         \~}
% \GetFileInfo{beamerthemeosx.dtx}
% \DoNotIndex{
%   \@empty, \@onlypreamble, \AtBeginDocument,
%   \newcounter, \addtocounter,
%   \newcommand,\renewcommand,\def,\edef,\let,\if,\else,\fi,
%   \begin, \begingroup, \endgroup, \end, \bfseries,
%   \centering, \clip, \copyright, \definecolor, \draw, \empty,
%   \footheight, \headheight, \headsep, \hfill, \href, \hspace,
%   \IfNoValueTF, \ifnum, \ifx, \includegraphics, \NewDocumentCommand,
%   \newif, \newlength, \normalsize, \parbox, \path, \ProvidesPackage,
%   \relax, \RequirePackage, \setcounter, \setlegnth, \small, \the, \tiny,
%   \vfill, \voffset, \vspace
% }
% \MakeShortVerb{\|}
% \NewDocElement[idxtype=color,printtype=color,idxgroup=color]
%   {Colour}{colour}
% \NewDocElement[macrolike=true,printtype=length]{Length}{length}
%
%% \makeatletter
% \def\ps@mystyle{
%   \def\@oddfoot{\hfil\thepage\hfil}
%   \def\@evenfoot{\hfil\thepage\hfil}
%   \def\@evenhead{\hfil\slshape\leftmark}
%   \def\@oddhead{\slshape\rightmark\hfil}}
% \makeatother
% \pagestyle{mystyle}
% \newcommand\btosxdocmarkright{%
%   \markright{\hskip-.5\marginparwidth \textsf{Beamer Theme OSX}
%     2026/09/30 v1.0}}
% \btosxdocmarkright
%
% \NewDocumentCommand{\btosxfoothref}{m m}{%
%   \href{#1}{#2}\footnote{\texttt{#1}}}
%
% \title{\textsf{Beamer Theme OSX\footnote{Copyright \copyright{} 2026 Didier
%   Verna. Released under the LPPL.}}\\
%   \itshape\normalsize A MacOS X 10.5 (Leopard) look for your slides}
% \author{Didier Verna\\
%   \texttt{mailto:didier@didierverna.net}\\
%   \texttt{https://didierverna.net}}
% \date{2026/09/30 v1.0}
%
% \maketitle
%
% \begin{center}
%   \includegraphics[width=.4\linewidth]{btosxlogo}
% \end{center}
%
% \tableofcontents
%
% \section{Installation}
% \textbf{Resources}
% \begin{itemize}\setlength\itemsep{0pt}
% \item \btosxfoothref{didierverna.net/projects/doceng/beamer-theme-osx}{%
%   Project homepage}
% \item \btosxfoothref{github.com/didierverna/beamer-theme-osx}{%
%   Project repository}
% \end{itemize}
% From the source directory, type \texttt{l3build install} to install the
% theme locally. Type \texttt{l3build doc} to compile the documentation. After
% the theme is installed, you can also switch to the \texttt{demo/} directory
% and compile the example there. This will allow you to see what the theme
% looks like. The demonstration source file also contains all the
% customization commands for quick reference.
%
% \smallskip
%
% From a CTAN archive, run the file \texttt{beamerthemeosx.ins} through
% \LaTeX. Then, copy all style and image files to your TDS-compliant
% installation, typically: \texttt{[TEXMF]/tex/latex/beamer-theme-osx/}. The
% documentation is already compiled. You may also process the demonstration
% file \texttt{demo.tex}. This will allow you to see what the theme looks
% like. The demonstration source file's preamble contains all the
% customization commands for quick reference.
%
% \smallskip
%
% From a TDS-compliant archive (\texttt{.tds.zip}), unpack directly into your
% \texttt{[TEXMF]} directory. The precompiled documentation and demo files
% mentioned above will be located in
% \texttt{[TEXMF]/doc/latex/beamer-theme-osx/}.
%
% \section{Usage}
% Load the theme in the usual Beamer way:
% \begin{verbatim}
% \usetheme{osx}
% \end{verbatim}
%
% \subsection{Social network links}
% In addition to its usual contents, the title page can display links to
% several social networks. The links appear as circular logos with a label
% below. The currently supported links are the following (they appear in that
% order): personal website, X, Facebook, and LinkedIn. Here is how to specify
% those links.
%
% \smallskip
%
% \noindent\DescribeMacro{\btosxwebsite}\oarg{label}\marg{url}\\
% Specifiy a personal website. The label defaults to the url itself. If you
% use this macro, you must also include an image file named
% \texttt{btosxwebsite} in your source directory.\\
% \DescribeMacro{\btosxx}\oarg{label}\marg{account}\\
% Specify an X account. The label defaults to \texttt{@account}.\\
% \DescribeMacro{\btosxfacebook}\oarg{label}\marg{account}\\
% Specify a Facebook account. The label defaults to \texttt{account}.\\
% \DescribeMacro{\btosxlinkedin}\oarg{label}\marg{account}\\
% Specify a LinkedIn account. The label defaults to \texttt{in/account}.
%
% \smallskip
%
% One caveat: beware of slashes in labels (or, more generally, characters with
% a depth) as they have a tendency to mess up the vertical alignment. If the
% problem occurs, simply put those characters in a \cs{raisebox} to cancel
% their depth.
%
% \subsection{Logo}
% \DescribeMacro{\logo}
% By default, the theme's logo appears at the top left corner of every slide.
% If you prefer no logo at all, or a logo of your own, simply use Beamer's
% regular \cs{logo} mechanism.
%
% \subsection{Copyright}
% \DescribeMacro{\btosxcopyright}\oarg{years=,name=,string=}\\
% The theme can display a copyright statement at the bottom left of each
% slide. The statement is of the form ``Copyright \copyright{} \meta{years}
% \meta{name}, \meta{string}''. All arguments can be empty, but the statement
% will only be typeset if at least a name or some years are provided.
% \meta{years} defaults to |\the\year|, \meta{name} defaults the Beamer's
% notion of a short author's name, and \meta{string} defaults to nothing.
% Example:
% \begin{verbatim}
% \btosxcopyright[years=2000--2026, name=My Lab, string=all rights reserved.]
% \end{verbatim}
%
% \subsection{Build information}
% \DescribeMacro{\btosxbuild}\marg{string}\\
% The theme can display a build string at the bottom right of each slide. This
% is useful if you want to advertise a build date, repository tag, or
% whatever.
%
% \subsection{Navigation bar}
% In lieu of a navigation bar, the theme displays each section as a desktop
% folder, in a horizontal row at the top right of each slide. This is the only
% sectionning level available. The current section is highlighted thanks to a
% slightly different icon, and a colored label.
%
% \smallskip
%
% \noindent\DescribeColour{BTOSXCurrentSectionColor} The color used to print
% the current section's (short) name under the folder icon. The default is set
% to yellow.
%
% \smallskip
%
% \noindent\DescribeLength{\BTOSXFolderSep} The length used to separate each
% section folder's icon in the navigation area. The default is 30pt, which may
% occasionally be too small if you have long (short) section names.
%
% \MaybeStop{\PrintChanges}
%
% \section{Implementation}
% \changes{v1.0}{2026/09/28}{First public version}
%
% \subsection{The main theme}
%    \begin{macrocode}
%<*theme>
\ProvidesPackage{beamerthemeosx}[2026/09/30 v1.0 Beamer Theme OSX - Main]

%    \end{macrocode}
% \begin{macro}{\btosxwebsite}
%    \begin{macrocode}
\@onlypreamble\btosxwebsite
\NewDocumentCommand{\btosxwebsite}{o m}{
  \def\btosx@websiteurl{#2}
  \def\btosx@websitelabel{\IfNoValueTF{#1}{#2}{#1}}}
\btosxwebsite{}

%    \end{macrocode}
% \end{macro}
% \begin{macro}{\btosxx}
%    \begin{macrocode}
\@onlypreamble\btosxx
\NewDocumentCommand{\btosxx}{o m}{
  \def\btosx@xurl{#2}
  \def\btosx@xlabel{\IfNoValueTF{#1}{@#2}{#1}}}
\btosxx{}

%    \end{macrocode}
% \end{macro}
% \begin{macro}{\btosxfacebook}
%    \begin{macrocode}
\@onlypreamble\btosxfacebook
\NewDocumentCommand{\btosxfacebook}{o m}{
  \def\btosx@facebookurl{#2}
  \def\btosx@facebooklabel{\IfNoValueTF{#1}{#2}{#1}}}
\btosxfacebook{}

%    \end{macrocode}
% \end{macro}
% \begin{macro}{\btosxlinkedin}
%    \begin{macrocode}
\@onlypreamble\btosxlinkedin
\NewDocumentCommand{\btosxlinkedin}{o m}{
  \def\btosx@linkedinurl{#2}
  \def\btosx@linkedinlabel{%
    %% The slash messes up vertical alignment.
    \IfNoValueTF{#1}{in\raisebox{0pt}[0pt][0pt]{/}#2}{#1}}}
\btosxlinkedin{}

%    \end{macrocode}
% \end{macro}
% \begin{macro}{\btosxcopyright}
%    \begin{macrocode}
\pgfkeys{
  /btosx/copyright/years/.store in  = \btosx@copyrightyears,
  /btosx/copyright/name/.store in   = \btosx@copyrightname,
  /btosx/copyright/string/.store in = \btosx@copyrightstring}

\newif\ifbtosx@copyright
\@onlypreamble\btosxcopyright
\NewDocumentCommand{\btosxcopyright}{O{}}{
  \def\btosx@copyrightyears{\the\year}
  \def\btosx@copyrightname{\beamer@shortauthor}
  \def\btosx@copyrightstring{}
  \pgfqkeys{/btosx/copyright}{#1}
  \btosx@copyrightfalse
  \ifx\btosx@copyrightyears\empty\else\btosx@copyrighttrue\fi
  \ifx\btosx@copyrightname\empty\else\btosx@copyrighttrue\fi
}
\btosxcopyright[name=,years=]

%    \end{macrocode}
% \end{macro}
% \begin{macro}{\btosxbuild}
%    \begin{macrocode}
\newcommand*\btosxbuild[1]{\def\btosx@build{#1}}
\btosxbuild{}

%    \end{macrocode}
% \end{macro}
%    \begin{macrocode}
\pgfdeclareimage[height=20pt]{btosxlogo}{btosxlogo}
\logo{\pgfuseimage{btosxlogo}}

\usefonttheme{osx}
\usecolortheme{osx}
\useinnertheme{osx}
\useoutertheme{osx}

\mode<presentation>

\RequirePackage{tikz}
\usetikzlibrary{calc}
\usetikzlibrary{patterns}
\usetikzlibrary{shadings}
\usetikzlibrary{shadows}

\pgfdeclareimage[width=\paperwidth,height=\paperheight]{%
  btosxbackground}{btosxbackground}
\setbeamertemplate{background canvas}{\pgfuseimage{btosxbackground}}

\setbeamersize{text margin left=16pt,text margin right=16pt}
\newlength\btosx@windowleft
\newlength\btosx@windowright
\newlength\btosx@windowtop
\newlength\btosx@windowbottom
\newlength\btosx@texttop
\newlength\btosx@textbottom
%% #### NOTE: we can't use \AtBeginDocument here because Beamer messes up with
%% the document environment. This, however, is Beamer's internal view of it.
\g@addto@macro\beamer@lastminutepatches{%
  \setlength\btosx@windowleft{\Gm@lmargin - .5\beamer@leftmargin}
  \setlength\btosx@windowright{\Gm@rmargin - .5\beamer@rightmargin}
  \setlength\btosx@windowtop{%
    %% #### NOTE: in fact, I could probably just use \headheight here because
    %% Beamer arranges for everything else to sum up to 0. However, let's not
    %% assume the dimensions wouldn't be changed outside Beamer itself.
    1in + \voffset + \topmargin + \headheight + \headsep}
  \setlength\btosx@windowbottom{\footheight}
  %% #### NOTE: the geometry of the gray bars and window buttons is currently
  %% hard-wired here and below.
  \setlength\btosx@texttop{\btosx@windowtop + 14pt}
  \setlength\btosx@textbottom{\footheight + 10pt}}

\pgfdeclareimage[height=8pt]{btosxtopleftbuttons}{btosxtopleftbuttons}
\pgfdeclareimage[height=8pt]{btosxtoprightbutton}{btosxtoprightbutton}
\setbeamertemplate{background}{%
  \ifnum\thepage>1\relax%
    \begin{tikzpicture}[remember picture,overlay]
      \path[fill=white,pattern=horizontal lines light gray, rounded corners,
          drop shadow={fill=black,opacity=.5,
            shadow xshift=.3ex,shadow yshift=-.3ex}]
        ($(current page.north west) + (\btosx@windowleft, -\btosx@windowtop)$)
        rectangle
        ($(current page.south east) - (\btosx@windowright, -\btosx@windowbottom)$);
      \clip[rounded corners]
        ($(current page.north west) + (\btosx@windowleft, -\btosx@windowtop)$)
        rectangle
        ($(current page.south east) - (\btosx@windowright, -\btosx@windowbottom)$);
      \path[shade,top color=btosx@toptopgray,bottom color=btosx@topbottomgray]
        ($(current page.north west) + (\btosx@windowleft, -\btosx@windowtop)$)
        rectangle
        ($(current page.north east) - (\btosx@windowright, \btosx@texttop)$);
      \node[anchor=north west,inner sep=0pt]
        at ($(current page.north west)
        + (\btosx@windowleft + 5pt, -\btosx@windowtop - 3pt)$)
        {\pgfuseimage{btosxtopleftbuttons}};
      \node[anchor=north east,inner sep=0pt]
        at ($(current page.north east)
        - (\btosx@windowright + 5pt, \btosx@windowtop + 3pt)$)
        {\pgfuseimage{btosxtoprightbutton}};
      \draw[color=btosx@grayline]
        ($(current page.north west) + (\btosx@windowleft, -\btosx@texttop)$)
        --
        ($(current page.north east) - (\btosx@windowright, +\btosx@texttop)$);
      \path[shade,top color=btosx@bottomtopgray,bottom color=btosx@bottombottomgray]
        ($(current page.south west) + (\btosx@windowleft, \btosx@textbottom)$)
        rectangle
        ($(current page.south east) - (\btosx@windowright, -\btosx@windowbottom)$);
      \draw[color=btosx@grayline]
        ($(current page.south west) + (\btosx@windowleft, \btosx@textbottom)$)
        --
        ($(current page.south east) - (\btosx@windowright, -\btosx@textbottom)$);
      \end{tikzpicture}
  \fi}

\setbeamertemplate{blocks}[rounded][shadow=true]

\mode<all>
%</theme>
%    \end{macrocode}
%
% \subsection{The inner theme}
%    \begin{macrocode}
%<*inner>
\ProvidesPackage{beamerinnerthemeosx}[2026/09/30 v1.0 Beamer Theme OSX - Inner]

\mode<presentation>

\RequirePackage{tikz}
\RequirePackage{xcolor}
\usetikzlibrary{positioning}
\usetikzlibrary{svg.path}

%% Official brand colours for the inline icons below.
\definecolor{btosx@facebookblue}{RGB}{24,119,242}
\definecolor{btosx@linkedinblue}{RGB}{10,102,194}

%% The X, Facebook and LinkedIn logos are drawn inline as real vector paths
%% Each path below has been pre-translated (in its own coordinate system) so
%% that it is centred on (0,0); this lets us scale it into place with a plain
%% xscale/yscale, sidestepping a couple of TikZ quirks (a parser bug on the
%% SVG "z" then relative "m" idiom, and "path picture" not playing well with
%% nested scaled/flipped scopes). The raw path data is inlined directly at
%% each point of use below (rather than behind a \newcommand) because TikZ's
%% "svg" path operation scans its argument itself and does not expand macros
%% in it.
\tikzset{
  title page link/.style={%
    circle,draw,thick,minimum width=.8cm,inner sep=0pt},
  title page link website/.style={%
    path picture={\node at (path picture bounding box.center) {%
        \includegraphics[width=.8cm]{btosxwebsite}};}},
  title page link label/.style={label=below:\parbox{1.5cm}{\centering#1}}}

\setbeamertemplate{title page}
{
  \begingroup
    \centering
    \vfill
    \begin{beamercolorbox}[center]{title}
      \usebeamerfont{title}\inserttitle
    \end{beamercolorbox}
    \ifx\insertsubtitle\@empty\else
      \begin{beamercolorbox}[center]{subtitle}
        \usebeamerfont{subtitle}\insertsubtitle
      \end{beamercolorbox}
    \fi
    \vspace{2ex}
    \begin{beamercolorbox}[center]{author}
      \usebeamerfont{author}\insertauthor
    \end{beamercolorbox}
    \begin{beamercolorbox}[center]{institute}
      \usebeamerfont{institute}\insertinstitute
    \end{beamercolorbox}
    \vspace{2ex}
    \begin{beamercolorbox}[center]{date}
      \usebeamerfont{date}\insertdate
    \end{beamercolorbox}
    \vspace{5ex}
    \usebeamercolor{links}\usebeamerfont{links}
    \begin{tikzpicture}[color=fg,fill=bg]
      \ifx\btosx@websiteurl\@empty\else
        \node[title page link,title page link website,
          title page link label={\href{\btosx@websiteurl}{\btosx@websitelabel}}]
          {};
      \fi
    \end{tikzpicture}
    \begin{tikzpicture}[color=fg,fill=bg]
      \ifx\btosx@xurl\@empty\else
        \node[title page link,fill=black,
          title page link label={%
            \href{http://www.x.com/\btosx@xurl}{\btosx@xlabel}}]
          {};
        \begin{scope}[xscale=.5cm/24,yscale=-.5cm/24]
          \path[fill=white] svg{%
            M 6.244 -9.75h 3.308l -7.227 8.26 8.502 11.24H 4.17l -5.214
            -6.817L -7.01 9.75H -10.32l 7.73 -8.835L -10.746 -9.75H
            -3.92l 4.713 6.231z
            M 5.083 7.77h 1.833L -4.916 -7.874H -6.883z};
        \end{scope}
      \fi
    \end{tikzpicture}
    \begin{tikzpicture}[color=fg,fill=bg]
      \ifx\btosx@facebookurl\@empty\else
        \node[title page link,fill=btosx@facebookblue,
          title page link label={%
            \href{http://www.facebook.com/\btosx@facebookurl}{%
              \btosx@facebooklabel}}] {};
        \begin{scope}[xscale=.46cm/15.5465,yscale=-.46cm/15.5465]
          \path[fill=white] svg{%
            M 1.2802 0.6818v 7.0915H -1.755v -7.0915h -2.5381v -2.7637h
            2.5381v -2.0382c 0 -2.3652 1.5363 -3.6531 3.7801 -3.6531
            1.0749 0 1.9987 0.0752 2.268 0.1089v 2.4719l -1.5566
            0.0006c -1.2201 0 -1.4563 0.5453 -1.4563 1.3454v 1.7645h
            2.9104l -0.3789 2.7637h -2.5315Z};
        \end{scope}
      \fi
    \end{tikzpicture}
    \begin{tikzpicture}[color=fg,fill=bg]
      \ifx\btosx@linkedinurl\@empty\else%
        \node[title page link,fill=btosx@linkedinblue,
          title page link label={%
            \href{http://www.linkedin.com/in/\btosx@linkedinurl}{%
              \btosx@linkedinlabel}}] {};
        \begin{scope}[xscale=.4cm/10.708,yscale=-.4cm/10.708]
          \path[fill=white] svg{%
            M -3.047 5.354V -1.871H -5.448v 7.225z
            M -4.247 -2.858c 0.837 0 1.358 -0.554 1.358 -1.248 -0.015
            -0.709 -0.52 -1.248 -1.342 -1.248S -5.59 -4.814 -5.59
            -4.106c 0 0.694 0.521 1.248 1.327 1.248z
            M 0.661 5.354V 1.319c 0 -0.216 0.016 -0.432 0.08 -0.586
            0.173 -0.431 0.568 -0.878 1.232 -0.878 0.869 0 1.216
            0.662 1.216 1.634v 3.865h 2.401V 1.21c 0 -2.22 -1.184
            -3.252 -2.764 -3.252 -1.274 0 -1.845 0.7 -2.165 1.193v
            0.025h -0.016l 0.016 -0.025V -1.871h -2.4c 0.03 0.678 0
            7.225 0 7.225z};
        \end{scope}
      \fi
    \end{tikzpicture}
    \vfill
  \endgroup
}

\mode<all>
%</inner>
%    \end{macrocode}
%
% \subsection{The outer theme}
%    \begin{macrocode}
%<*outer>
\ProvidesPackage{beamerouterthemeosx}[2026/09/30 v1.0 Beamer Theme OSX - Outer]

\mode<presentation>

\RequirePackage{calc}
\RequirePackage{tikz}
\usetikzlibrary{calc}

\setbeamertemplate{frametitle}{%
  \begin{tikzpicture}[remember picture,overlay,
    every node/.style={anchor=north west,align=left,inner sep=1pt}]
    \node at ($(current page.north west)
              + (\btosx@windowleft + 38pt,-\btosx@windowtop)$)
      {\insertframetitle};
  \end{tikzpicture}%
}


%% Logo
\setbeamertemplate{insert logo}{\tikz[inner sep=0pt] \node { \insertlogo };}


%% Headline
\newcounter{@btosxfolderposition}
%    \end{macrocode}
% \begin{length}{\BTOSXFolderSep}
%    \begin{macrocode}
\newlength\BTOSXFolderSep\setlength\BTOSXFolderSep{30pt}
%    \end{macrocode}
% \end{length}
%    \begin{macrocode}
%% #### NOTE: ideally, I would have preferred to tweak \voffset or \topmargin
%% here, but I cannot because because this would shift the whole Beamer frame
%% (including background pictures). So let's make this vertical space part of
%% the headline directly.
\newlength\@btosxheadlinesep\setlength\@btosxheadlinesep{4pt}

%% #### NOTE: Beamer determines the headline's height automatically, right at
%% the beginning of the document, as described in section 8.2.1 of the
%% manual. This means that I need to typeset the headline even on the title
%% page. Otherwise, Beamer would think it's size is 0.
\setbeamertemplate{headline}{%
  \setcounter{@btosxfolderposition}{0}%
  \vspace{\@btosxheadlinesep}
  \begin{beamercolorbox}{headline}%
    \insertsectionnavigationhorizontal{\paperwidth}{%
      %% The 4pt are because of the window's round corners. It looks better
      %% like this IMO.
      \hspace{\btosx@windowleft}\hspace{4pt}\usebeamertemplate{insert logo}%
      \hfill%
      \begin{tikzpicture}[inner sep=0pt,x=\BTOSXFolderSep]}{%
      \end{tikzpicture}%
      %% The 4pt are because of the window's round corners. It looks better
      %% like this IMO.
      \hspace{\Gm@rmargin}\hspace{4pt}}%
  \end{beamercolorbox}%
  \vspace{\@btosxheadlinesep}
}

%% #### NOTE: we need to use the font option in order to process Beamer color
%% names. The color option only understands xcolor names.

\pgfdeclareimage[height=15pt]{btosxfolder}{btosxfolder}
\setbeamertemplate{section in head/foot shaded}{%
  \node[label={[font={\usebeamercolor[fg]{section in head/foot shaded}}]%
      below:\insertsectionhead}]
    at (\the@btosxfolderposition,0)
    {\pgfuseimage{btosxfolder}};
  \addtocounter{@btosxfolderposition}{1}}

\pgfdeclareimage[height=15pt]{btosxfoldercurrent}{btosxfoldercurrent}
\setbeamertemplate{section in head/foot}{%
  \node[label={[font={\usebeamercolor[fg]{section in head/foot}}]%
      below:\insertsectionhead}]
    at (\the@btosxfolderposition,0)
    {\pgfuseimage{btosxfoldercurrent}};
  \addtocounter{@btosxfolderposition}{1}}


%% Footline

\setbeamertemplate{window footline left}{%
  \insertshorttitle%
  \ifx\beamer@shortsubtitle\empty\else~/~\insertshortsubtitle\fi}
\setbeamertemplate{window footline center}{}
\setbeamertemplate{window footline right}{%
  \insertframenumber/\inserttotalframenumber}

\setbeamertemplate{window footline}{%
  \ifnum\thepage>1%
    \begin{tikzpicture}[overlay,remember picture,
      every node/.style={inner xsep=5pt,inner ysep=2pt}]
      \node[anchor=north west,align=left]
        at ($(current page.south west) + (\btosx@windowleft, \btosx@textbottom)$)
        { \usebeamertemplate***{window footline left} };
      \node[anchor=north west,align=center,text width=\paperwidth]
        at ($(current page.south west) + (\btosx@windowleft, \btosx@textbottom)$)
        { \usebeamertemplate***{window footline center} };
      %% #### WARNING: the 3pt inner xsep overrides the global 5pt one in
      %% order to compensate for a right margin bigger the the left one, for a
      %% reason I don't understand.
      \node[anchor=north east,align=right,inner xsep=3pt]
        at ($(current page.south east) - (\btosx@windowright, -\btosx@textbottom)$)
        { \usebeamertemplate***{window footline right} };
    \end{tikzpicture}%
  \fi%
}

\setbeamertemplate{footline}{%
  \usebeamertemplate***{window footline}
  \begin{beamercolorbox}[leftskip=2pt,rightskip=4pt]{navigation symbols}
    \raisebox{2pt}{%
      \ifbtosx@copyright%
        Copyright \copyright{} \btosx@copyrightyears{} \btosx@copyrightname%
        \ifx\btosx@copyrightstring\empty\else, \btosx@copyrightstring\fi%
      \fi}%
    \hfill%
    \usebeamertemplate***{navigation symbols}%
    \hfill%
    \raisebox{2pt}{\ifx\btosx@build\empty\else\btosx@build\fi}%
  \end{beamercolorbox}
  %% Beamer adds 4pt automatically before the footline.
  \vspace{2pt}}

%% Need this for now because the default one typesets the logo and the navbar
%% at the bottom right corner of the frame.
\setbeamertemplate{sidebar right}{}

\mode<all>
%</outer>
%    \end{macrocode}
%
% \subsection{The color theme}
%    \begin{macrocode}
%<*color>
\ProvidesPackage{beamercolorthemeosx}[2026/09/30 v1.0 Beamer Theme OSX - Color]

\mode<presentation>
%    \end{macrocode}
% \begin{colour}{BTOSXCurrentSectionColor}
%    \begin{macrocode}
\definecolor{BTOSXCurrentSectionColor}{named}{yellow}
%    \end{macrocode}
% \end{colour}
%    \begin{macrocode}
%% For the frame components:
\definecolor{btosx@bottombottomgray}{RGB}{147,147,147}
\definecolor{btosx@bottomtopgray}{RGB}{181,181,181}
\definecolor{btosx@topbottomgray}{RGB}{186,186,186}
\definecolor{btosx@toptopgray}{RGB}{198,198,198}
\definecolor{btosx@grayline}{RGB}{85,85,85}

%% For blocks:
\definecolor{btosx@blocktitlebackground}{RGB}{198,198,198}
\definecolor{btosx@blockbodybackground}{RGB}{226,230,232}

%% For various uses, including alerts:
\definecolor{btosx@blue}{RGB}{55,103,167}

\setbeamercolor{background canvas}{fg=white,bg=}

\setbeamercolor{title}{parent=background}
\setbeamercolor{subtitle}{parent=background}
\setbeamercolor{author}{parent=background}
\setbeamercolor{institute}{parent=background}
\setbeamercolor{date}{parent=background}
\setbeamercolor{links}{parent=background}

\setbeamercolor{headline}{parent=background}
\setbeamercolor{section in head/foot shaded}{parent=headline}
\setbeamercolor{section in head/foot}{parent=section in head/foot shaded,
  fg=BTOSXCurrentSectionColor}

\setbeamercolor{normal text}{fg=black,bg=white}

\setbeamercolor{frametitle}{parent=normal text}

\setbeamercolor{alerted text}{fg=btosx@blue}
\setbeamercolor{example text}{fg=btosx@blue}

\setbeamercolor{structure}{fg=btosx@blue}

\setbeamercolor{block title}{fg=black,bg=btosx@blocktitlebackground}
\setbeamercolor{block body}{bg=btosx@blockbodybackground}
\setbeamercolor{block title alerted}{bg=btosx@blocktitlebackground}
\setbeamercolor{block body alerted}{parent=block body}
\setbeamercolor{block body example}{parent=block body}
\setbeamercolor{block title example}{parent=block title}

\setbeamertemplate{section in toc shaded}[default][50]

\mode<all>
%</color>
%    \end{macrocode}
%
% \subsection{The font theme}
%    \begin{macrocode}
%<*font>
\ProvidesPackage{beamerfontthemeosx}[2026/09/30 v1.0 Beamer Theme OSX - Font]

\mode<presentation>

\setbeamerfont{title}{series=\bfseries}
\setbeamerfont{date}{size=\small}
\setbeamerfont{links}{parent=structure,size=\tiny}

\setbeamerfont{frametitle}{parent=normal text,series=\bfseries}
\setbeamerfont{framesubtitle}{size=\normalsize}

\setbeamerfont{headline}{series=\bfseries}
\setbeamerfont{section in head/foot shaded}{parent=headline}
\setbeamerfont{section in head/foot}{parent=section in head/foot shaded}

\setbeamerfont{block title}{series=\bfseries}

\mode<all>
%</font>
%    \end{macrocode}
%
% \Finale
\endinput
